\section{量化与蒸馏把成本换成精度}
\label{sec:14-Deep-Learning/09-Scaling-and-Adaptation/04}

模型训好了，评测漂亮，现在要上线。约束是硬的：单卡显存放不下，首 token 延迟必须压到某个毫秒数以内，而每天的调用量还在涨。重训一个小模型来不及，也不划算——那等于把上一节算过的账重付一遍。

能动的只剩两处。一是让同一个模型占更少的位：权重和激活从 \(16\) 位降到 \(8\) 位、\(4\) 位，显存与访存带宽按位宽比例下降。二是换一个更小的模型顶上，但让它去学大模型的输出而不是原始标签。这两条路听起来都是``用精度换成本''，实际上动的是完全不同的东西，混着谈会得出错误的结论。

\subsection{量化：改变数值表示，不改变函数}

均匀量化把一个实数区间切成等份。给定位宽 \(b\) 和取值范围 \([m, M]\)，步长是

\[
  \Delta = \frac{M - m}{2^{b} - 1},
\]

每个数被映射到最近的格点上。舍入误差 \(e\) 落在 \([-\Delta/2, \Delta/2]\) 内，若近似当作均匀分布处理，

\[
  \E[e] = 0, \qquad \Var(e) = \frac{\Delta^{2}}{12}.
\]

这个式子把整件事的经济性讲清楚了：每多一位，\(\Delta\) 减半，误差方差降到四分之一。反过来，从 \(8\) 位降到 \(4\) 位省一半显存，误差的标准差涨十六倍。省的是线性的，付的是指数的——和上一单元那条幂律恰好是同一种不对称，只是方向相反。

这里的``均匀分布''是一个方便的近似而非事实。它要求被量化的值相对步长足够分散且与格点无系统性对齐；权重分布通常大致满足，某些高度集中或稀疏的分布则不满足，此时真实误差会偏离这个方差。关于舍入误差如何在一串运算中产生和累积，\hyperref[sec:08-Numerical-Analysis/01-Floating-Point-and-Error/02]{``误差如何产生与放大''}给的是同一套账。

\subsection{离群值才是难点}

权重量化通常比预期顺利，原因在分布形状：训练好的权重大致集中在零附近，尾部不长，于是 \([m, M]\) 不必取得很宽，\(\Delta\) 就小，相对误差可控。

激活值不是这样。同一层的激活里常常存在少数几个维度，其数值比其余维度大一到两个数量级，而且位置相当稳定——同一批离群维度在不同输入上反复出现。麻烦在于 \(\Delta\) 由 \(M - m\) 决定，也就是由最大的那个离群值决定。一个大一百倍的离群值把 \(\Delta\) 撑大一百倍，剩下百分之九十九的值原本分布在一个窄区间里，现在全被挤进两三个格点，等于被抹平。位宽花在了几乎没有信息的动态范围上。

针对这一点的做法都在同一个方向上：不让一个缩放因子管太多数。分组量化按通道或按小块各算各的 \([m, M]\)，离群维度多大都只影响它自己那一组；也可以把识别出的离群维度整体留在高精度，其余部分低位宽处理。两种做法都在承认同一件事——量化的困难不在位宽，在于共享一个缩放因子的那组数是否同质。

误差还会累积。第 \(\ell\) 层引入的扰动 \(\delta_\ell\) 会被后面所有层继续变换，粗略地看，它被放大的倍数与后续各层线性映射的算子范数之积同阶。深度一大，这个乘积可以很大，靠前的层因此更敏感。实践里被反复观察到的是：第一层、最后一层、以及注意力里的部分投影，量化后退化明显大于其他层，通常保留较高精度。这个``敏感层''清单是经验清单，不同架构上并不完全一样，也没有一个原理性的判据能事先把它算出来——现有做法是逐层试着降位宽，看整体指标什么时候开始掉。

要点在于：量化没有换掉那个函数。它试图用更少的位表示同一组权重、算出同一组数，误差是逼近误差，理论上位宽给够就能任意接近。

\subsection{蒸馏：换掉函数，但换一个更好教的目标}

蒸馏做的是另一件事——训练一个参数少得多的学生模型，去拟合教师模型的输出。关键在于拟合什么。

如果只用原始的硬标签训小模型，每个样本提供的监督是一个类别编号，损失只关心正确类的那一个概率。用教师的完整输出分布代替它，同一个样本提供的是整个概率向量。这两者携带的信息量差别很大：教师说``这张图有 \(0.7\) 的把握是狼、\(0.2\) 是狗、\(0.001\) 是卡车''，不只告诉学生答案，还告诉了它狼与狗在这个模型眼里很近而与卡车很远。这份类别间的相似结构是教师从海量数据里提炼出来的，硬标签里完全没有。

把它写成目标函数就是散度最小化。记教师和学生在样本 \(x\) 上的输出分布为 \(P_T(\cdot \given x)\) 与 \(P_S(\cdot \given x)\)，蒸馏损失是

\[
  \mathcal{L}_{\text{distill}} = \E_{x}\left[\, \KL\bigl(P_T(\cdot \given x) \,\big\|\, P_S(\cdot \given x)\bigr) \right].
\]

因为教师固定，\(\KL\) 里教师的熵是常数，最小化它等价于最小化交叉熵——这正是\hyperref[sec:07-Information-Theory/03-Divergences/01]{``错误模型的代价：从交叉熵到 KL 散度''}给出的等价。硬标签训练是这个式子的退化情形：把 \(P_T\) 换成 one-hot 分布，那些非零的小概率连同它们承载的相似结构一起被丢掉。方向也值得留意，这里用的是 \(\KL(P_T \| P_S)\)，它惩罚学生在教师有概率的地方给出零概率，逼学生覆盖教师的全部支撑，而不是只抓住主模式。

温度参数处理的是另一个问题。教师的分布常常极尖，正确类的概率接近 \(1\)，其余类之间的相对大小虽然含有信息，但数值太小，在损失里的权重可以忽略。把 logits 除以温度 \(\tau > 1\) 再取 softmax，

\[
  P^{(\tau)}(y \given x) = \frac{\exp(z_y / \tau)}{\sum_{y'} \exp(z_{y'} / \tau)},
\]

分布被软化，小概率之间的比例关系被抬到可见的量级。教师和学生用同一个 \(\tau\)。软化的副作用是梯度整体按 \(1/\tau^{2}\) 缩小，所以这一项通常再乘回 \(\tau^{2}\)，让 \(\tau\) 的调整不干扰学习率。

蒸馏的效果同样是经验性的。它常常让小模型明显好过同尺寸从头训练的模型，但没有定理保证一定如此，也没有办法事先知道某个教师--学生尺寸比能保住多少能力。学生比教师小得越多，能被搬过去的越少，这条趋势可靠，具体的界不存在。

\subsection{它们不在同一条曲线上}

\begin{center}
\begin{tikzpicture}[scale=1.0]

\draw[->,>=stealth] (0.4,0.4) -- (9.2,0.4);
\draw[->,>=stealth] (0.4,0.4) -- (0.4,4.60);
\node[anchor=north] at (4.7,0.20) {\footnotesize 每 token 的显存与延迟};
\node[anchor=south,rotate=90] at (-0.05,2.50) {\footnotesize 质量};

% 原模型量化曲线
\draw[very thick] plot[smooth] coordinates {(7.60,3.80)(5.20,3.70)(3.30,3.25)(2.20,2.30)};
\fill (7.60,3.80) circle (2.2pt);
\fill (5.20,3.70) circle (2.2pt);
\fill (3.30,3.25) circle (2.2pt);
\fill (2.20,2.30) circle (2.2pt);
\node[anchor=south] at (7.60,3.92) {\scriptsize 原模型 \(16\) 位};
\node[anchor=south] at (5.20,3.82) {\scriptsize \(8\) 位};
\node[anchor=south] at (3.30,3.37) {\scriptsize \(4\) 位};
\node[anchor=east] at (2.05,2.28) {\scriptsize \(3\) 位};

% 学生模型量化曲线
\draw[very thick,dashed] plot[smooth] coordinates {(3.90,1.80)(2.60,1.65)(1.60,0.85)};
\fill (3.90,1.80) circle (2.2pt);
\fill (2.60,1.65) circle (2.2pt);
\fill (1.60,0.85) circle (2.2pt);
\node[anchor=north west] at (4.05,1.68) {\scriptsize 学生模型 \(16\) 位};
\node[anchor=north] at (2.60,1.53) {\scriptsize \(8\) 位};
\node[anchor=east] at (1.45,0.83) {\scriptsize \(4\) 位};

% 量化：沿曲线左滑
\node[anchor=west] at (0.65,4.10) {\scriptsize 量化：沿同一条曲线左滑};
\draw[->,>=stealth] (4.65,4.10) -- (3.95,4.10);

% 蒸馏：跳到另一条曲线
\draw[->,>=stealth,very thick] (7.20,3.45) -- (4.15,1.98);
\node[anchor=north west] at (5.90,2.45) {\scriptsize 蒸馏：跳到另一条曲线};

\end{tikzpicture}
\end{center}

\emph{实线是原模型在不同位宽下的位置，虚线是蒸馏出的学生模型的同一条曲线。量化让你沿着一条曲线往左滑，滑到某个位宽之后质量急转直下；蒸馏是整体跳到另一条曲线上，起点更靠左但也更低。要记住的是两者可以叠加而不是二选一——学生模型自己也能被量化，图上那条虚线就是证据；但两次损失也会叠加，且不一定只是相加。}

图里那两条曲线的转折位置和落差都是示意的。真实的位置随架构、任务和量化方案变化很大，唯一稳定的是形状：高位宽段近乎水平，越过某个位宽后急剧下降；学生曲线整体在教师曲线的左下方。

\subsection{评测集上没掉，不等于没退化}

量化与蒸馏都有一个共同的、被反复低估的风险：它们在标准评测上的损失往往很小——小到报告里可以写``基本无损''——而在评测覆盖不到的地方退化明显更大。

有几类地方特别容易出问题。长尾输入：罕见语言、生僻术语、异常格式，它们在评测集里样本极少，模型在这些输入上的表现由权重里那些小幅度、被量化优先牺牲掉的成分决定。多步推理：如果一次回答要连对若干步，每步的正确率下降 \(\varepsilon\)，整条链的成功率大致按 \((1 - \varepsilon)^{k}\) 衰减，单步上看不出的退化在 \(k\) 步之后是显著的。以及安全相关的行为：拒答边界、对诱导性提问的稳健性，这些是在微调阶段用相对少量数据塑造出来的，对权重扰动的抵抗力本来就弱，而它们在通用评测里几乎没有权重。

结论不是``不要量化''，而是评测的设计必须和压缩的方式匹配。压缩改变了模型，所以压缩之后被评的是一个新模型，用原来那套指标去确认它``没退化''，逻辑上只证明了它在那套指标上没退化。这和\hyperref[sec:11-Mathematical-Statistics/06-Resampling-and-Model-Selection/03]{``交叉验证评估的是完整训练流程''}是同一条判断在另一个位置上的应用：被评估的对象是整条流程的输出，流程改了一步，评估就得重做，而且得针对改动最可能伤到的地方专门补测。

\medskip

这一单元的四节走完了一条线。规模与损失的关系是幂律，指数很小，所以能力的每一次抬升都要按 \(2^{1/\alpha}\) 的倍数付钱；给定预算怎么分配，是一个有解析解的约束优化问题，而解取决于你把哪些成本写进了目标；训练完之后，低秩增量让改变行为的成本降到基座的百分之一量级；量化与蒸馏则决定这份能力能以什么代价送到用户手里。

贯穿其中的是同一件事：规模决定了能力的上限，适配与压缩决定了这份能力的交付价格，而这两侧的所有可靠结论几乎都是经验规律。幂律没有证明，最优配比依赖拟合出来的指数，低秩假设是观察，量化的容忍度是试出来的。能当定理用的只有那几条：\(L_\infty\) 是熵给的地板，Lagrange 条件在给定假设下给出唯一最优，Eckart--Young 界定了秩 \(r\) 能装下多少。把这三条和其余的经验规律分清楚，是在这个领域里做判断的前提。
